% TOP_PROBLEM: {"id": 20002526, "title": "A quantified Gibbons-Hawking N-body cluster model", "category_id": 16, "category": {"id": 16, "name": "physics", "display_name": "Mathematical Physics", "description": "Problems at the intersection of mathematics and physics.", "slug": "physics", "order_index": 16, "created_at": "2026-07-31T15:26:25.671Z"}, "status": "partially_solved", "statusQualification": "This is a methodological first step rather than a single terminal conjecture. It is related to 20002525 and 20002527 but requests a different approach; no particular Gibbons–Hawking cluster model is invented as the full source target. The catalogue title/status was paired with a machine-generated research attempt; that attempt is not peer-reviewed evidence of settlement. The statement covers the original source questions rather than only the title’s subcase."}
% FIRST_POSTED: 2026-10-11
% ENTRY_KIND: statement_only
% SUBMISSION: {"kind": "statement_only", "category": "top_problems", "problem_id": "20002526", "proposal_type": "established", "source": "American Institute of Mathematics, $L^2$ harmonic forms in geometry and string theory; notes by Anda Degeratu and Mark Haskins, March 2004, 21 March 2004, Hausel wrap-up programme and first step 2. https://aimath.org/WWN/l2harmonic/l2harmonic.pdf", "status_evidence": "Inherited source-catalogue status partially_solved preserved. Background has machine-generated research claims explicitly marked not peer reviewed; these are not independent status evidence. Primary source accessed 11 October 2026; compared against the supplied catalogue excerpt. Historical wording and partial-result remarks were checked for scope; no comprehensive current-literature or proof-verification audit was performed.", "status_checked_on": "2026-10-11", "duplicate_check": "Full 15,773-record source registry checked for identical normalized statements; available site registry checked for identical normalized title/statement and canonical-ID collision. Canonical ID is absent from the site registry. Related canonical records 20002525, 20002527. No mathematical-equivalence literature audit performed.", "attribution": "Rafe Mazzeo, within Tamás Hausel’s programme. Statement preparation: OpenAI GPT-6 Codex assistant, generation 11 October 2026; source-matching and notation editing by the assistant, no independent human mathematical verification."}
% SOURCE_STATEMENT: "2. (Mazzeo) Combine asymptotics with N -body techniques."
% Preparation and final self-audit: CONTRIBUTING.md prompts C and D; source/audit artifacts retained in the local physics_statements_2026_10_11 batch.
% Imported from: physics_statements.tex; UnsolvedMath ID 20002526
% !TeX program = lualatex
\documentclass[11pt]{article}
\usepackage[margin=25mm]{geometry}
\usepackage{fontspec}
\setmainfont{Latin Modern Roman}
\usepackage{amsmath,amssymb,mathtools}
\usepackage{unicode-math}
\setmathfont{Latin Modern Math}
\usepackage{xurl,hyperref}
\hypersetup{colorlinks=true,urlcolor=blue}
\setcounter{secnumdepth}{0}
\setlength{\parindent}{0pt}
\setlength{\parskip}{5pt}
\setlength{\emergencystretch}{3em}
\newcommand{\sref}[2]{\hyperlink{source:#1:#2}{[S#2]}}
\newcommand{\cataloguescope}[1]{\paragraph{Recorded target.}#1}
\newcommand{\R}{\mathbb R}
\newcommand{\C}{\mathbb C}
\newcommand{\N}{\mathbb N}
\newcommand{\Z}{\mathbb Z}
\newcommand{\Q}{\mathbb Q}
\newcommand{\F}{\mathbb F}
\newcommand{\eps}{\varepsilon}
\begin{document}
% UnsolvedMath ID 20002526; source catalogue code AIM-PHYSICS-0034
\setcounter{section}{20002525}
\section{A quantified Gibbons-Hawking N-body cluster model}\label{20002526}
{\small UnsolvedMath ID 20002526 · Mathematical Physics}
\subsection{Definitions and mathematical statement}

\paragraph{Shared notation.}$\mathbb N=\{1,2,\ldots\}$; the standard number fields have their usual meanings. All additional conventions are stated locally.


\paragraph{Definitions and notation.}A hyper-Kähler metric has three parallel complex structures satisfying the quaternionic relations. In four dimensions a gravitational instanton is a complete noncompact Ricci-flat geometry in the source’s self-dual setting. ALE and ALF describe different asymptotic geometries; ALF has an asymptotic circle direction over a three-dimensional base. Toric hyper-Kähler quotients are hyper-Kähler reductions of flat quaternionic spaces by tori. N-body analytical methods organize asymptotic regions by clusters separating from one another and analyze operators or geometry in those regimes. Their application here requires compatible cluster decompositions and quantitative estimates on the moduli-space ends.

\paragraph{Catalogue wording retained for traceability.}
For the toric hyper-Kähler quotients in the source’s Hodge-cohomology programme, combine their asymptotic geometry with N-body techniques, as proposed by Mazzeo.
\paragraph{Scope and source qualification.}This is a methodological first step rather than a single terminal conjecture. It is related to 20002525 and 20002527 but requests a different approach; no particular Gibbons–Hawking cluster model is invented as the full source target. The catalogue title/status was paired with a machine-generated research attempt; that attempt is not peer-reviewed evidence of settlement. The statement covers the original source questions rather than only the title’s subcase.
\subsection{Short English statement}
Combine asymptotic geometry of toric hyper-Kähler quotients with N-body analytical techniques.
\subsection{Sources}
\begin{itemize}
\item[\hypertarget{source:20002526:1}{S1}] ULAM.AI and the UnsolvedMath Contributors, UnsolvedMath, record 20002526 (AIM-PHYSICS-0034). Catalogue curation under CC BY 4.0; definitions and English exposition added locally.
\url{https://huggingface.co/datasets/ulamai/UnsolvedMath}.
\item[\hypertarget{source:20002526:2}{S2}] American Institute of Mathematics, $L^2$ harmonic forms in geometry and string theory; notes by Anda Degeratu and Mark Haskins, March 2004, 21 March 2004, Hausel wrap-up programme and first step 2.
\url{https://aimath.org/WWN/l2harmonic/l2harmonic.pdf}.
\end{itemize}
\label{end:20002526}
\end{document}
